\begin{textsample}{POS Dim 2 – vem\_ed\_35 – Score 2.00 – vem\_ed\_35/t186.txt}  \label{ex:f2_pos_055}
Aos Leitores Osvaldo Succi Jr. - Coordenador de Projetos Este número de VEm destaca o papel transformador dos Projetos Colaborativos Internacionais (PCIs) na construção de uma educação superior mais conectada, inclusiva e global.
A quinta edição do International Virtual Exchange Symposium (IVES CGESG 2026), organizada pela equipe de Apoio à Internacionalização da Divisão de Extensão e Pesquisa no Ensino Superior da Coordenadoria Geral de Ensino Superior (CGESG) do Centro Paula Souza, reuniu especialistas de diferentes países para discutir desafios e caminhos da cooperação em Intercâmbios Virtuais.
Os palestrantes Sara Pittarello (UNICollaboration, rede europeia de apoio aos Intercâmbios Virtuais), Carlos Maza (UNAM, México), Henry Shepherd (IVEC, EUA) e Salome Escalona (Bukidnon State University, Filipinas) ressaltaram que a internacionalização vai além da tecnologia, demandando interação significativa, objetivos de aprendizagem compartilhados e parcerias sustentáveis.
A seção “Boas Práticas” apresenta duas experiências que revelam a dimensão humana desses projetos: PCIs com o Chile \textbf{conduzidos} pelo professor Alessandro Padin em 2025 revelaram como café e chocolate carregam identidade cultural.
PCIs realizados desde 2022 envolvendo diversas Fatecs e a Setsunan University (Japão) conquistaram escalabilidade a partir de um projeto com tarefas claras e bem-estruturadas, plataformas tecnológicas simples e amigáveis, diálogo constante entre os docentes e aperfeiçoamento contínuo ao longo dos \textbf{semestres}.
Destaca-se ainda a visita de alunos e professores da Kirkwood Community College (EUA) à Fatec Itaquera, complementando com atividades presenciais os PCIs que vêm sendo realizados virtualmente desde 2022.
O “Artigo de Opinião” conecta a temática do 5º IVES CGESG 2026 com a experiência do PCI realizado entre Fatecs Pompeia, Zona Leste e Guarulhos com Bukidnon State University Filipinas).
Boa leitura!

% matched lemmas: conduzir, semestr
\end{textsample}
